\subsection{有界测度的特征}

命题 12. --- 为使测度 \(\mu\) （局部紧空间 \(X\) 上的）是有界的（Ch. III, §1, No.~8），必须且只须 \(X\) 是关于 \(\mu\) 的可积集（或者，换一种说法，每个有限常值函数都可积）；此时

\[
\| \mu \| = | \mu | (\mathrm{X}) = \int d | \mu |.
\]

因为，我们已看到 \(|\mu|^{*}(X) = \|\mu\|\) （§1, No.~2）；于是命题由 No.~6 的命题 10 得出。

对每个有界测度 \(\mu\) ，我们仍称 \(\mu(\mathrm{X})\) 为 \(\mu\) 的总质量。

由 No.~6 的定理 4 可得：若 \(\mu\) 是一个有界测度，则对每个 \(\varepsilon > 0\) ，存在紧集 \(K\) 使 \(|\mu| (\mathbb{C}K) \leqslant \varepsilon\) 。

命题 13. --- 设 \(\mu\) 是 X 上的一个有界测度。设 \(\mathfrak{B}\) 是 \(\mathcal{L}_{\mathrm{F}}^{p}\) 上的一个滤子基，具有下列性质：

\(1^{\circ}\) 存在一个集 \(\mathbf{M} \in \mathfrak{B}\) ，使诸函数 \(\mathbf{f} \in \mathbf{M}\) 在 \(\mathbf{X}\) 上一致有界；

\(2^{\circ}\) \(\mathfrak{B}\) 在 X 的每个紧子集上一致收敛于一个函数 \(\mathbf{f}_0\) 。

在这些条件下，\(\mathbf{f}_0\) 属于 \(\mathcal{L}_{\mathrm{F}}^{p}\) ，且 \(\mathfrak{B}\) \(p\) 阶平均收敛于 \(\mathbf{f}_0\) 。

我们首先注意，若 \(|\mathbf{f}(x)| \leqslant a\) 对每个 \(x \in X\) 与每个函数 \(\mathbf{f} \in M\) 成立，则也有 \(|\mathbf{f}_0(x)| \leqslant a\) 对每个 \(x \in X\) 成立。这样，对每个 \(\varepsilon > 0\) ，存在紧集 \(K\) 使 \(|\mu|(\mathbb{C}K) \leqslant \varepsilon^p\) ，且存在集 \(N \in \mathfrak{B}\) 使对每个函数 \(\mathbf{f} \in N\) ，有 \(|\mathbf{f}(x) - \mathbf{f}_0(x)| \leqslant \varepsilon(|\mu|(K))^{-1/p}\) 对所有 \(x \in K\) 成立。现在，我们可以写

\[
\mathbf {f} - \mathbf {f} _ {0} = (\mathbf {f} - \mathbf {f} _ {0}) \varphi_ {\mathrm{K}} + (\mathbf {f} - \mathbf {f} _ {0}) \varphi_ {\mathbf {C K}};
\]

由前所述，若 \(\mathbf{f} \in \mathrm{M} \cap \mathrm{N}\) ，则 \(\mathrm{N}_p((\mathbf{f} - \mathbf{f}_0)\varphi_{\mathrm{K}}) \leqslant \varepsilon\) 且 \(\mathrm{N}_p((\mathbf{f} - \mathbf{f}_0)\varphi_{\mathbf{C}_{\mathrm{K}}}) \leqslant 2a\varepsilon\) ，由此 \(\mathrm{N}_p(\mathbf{f} - \mathbf{f}_0) \leqslant (2a + 1)\varepsilon\) ，这就证明了命题。

推论. --- 对有界测度 \(\mu\) （\(X\) 上的），每个有界连续映射 \(\mathbf{f}\) （\(X\) 到 \(\mathbf{F}\) 中的）都属于每个 \(\mathcal{L}_{\mathrm{F}}^{p}\) （ \(1 \leqslant p < +\infty\) ）。

对 X 的每个紧子集 K，设 \(M_{K}\) 是形如 hf 的 X 到 F 中的映射所成之集，其中 h 是 X 到 {[}0,1{]} 中的一个连续映射，它在 K 上等于 1 且具有紧支集。显然诸集 \(M_{K}\) 构成 \(L_{F}^{p}\) 上的一个滤子基 B，属于 \(M_{K}\) 的函数是一致有界的，且 B 在 X 的每个紧子集上一致收敛于 f，由此得推论。

特别地，函数 f 是可积的，且它的积分 \(\int f d\mu\) 是诸积分 \(\int h f d\mu\) 关于 B 的极限。

我们将在 §5, No.~6 中作为可积性的一般判别法的推论重新得到命题 13 的推论。

以 Ch. III, §1, No.~2 的记号，\(|\mathbf{f}| \leqslant \| \mathbf{f} \| \cdot 1\) 对每个函数 \(\mathbf{f} \in \mathcal{C}^b(\mathrm{X};\mathrm{F})\) 成立，于是由 No.~2 的公式 (3) 与 (4)，

\[
\mathrm{N} _ {p} (\mathbf {f}) \leqslant \| \mathbf {f} \| \cdot \mathrm{N} _ {p} (1) = \| \mathbf {f} \| \cdot \| \mu \| ^ {1 / p}.\tag{19}
\]

特别地，对 \(p = 1\) ，No.~2 的公式 (5) 给出

\[
\left| \int \mathbf {f} d \mu \right| \leqslant \| \mathbf {f} \| \cdot \| \mu \|,\tag{20}
\]

因而映射 \(f \mapsto \int f d\mu\) 在 Banach 空间 \(\mathcal{C}^{b}(X;F)\) 上是连续的；它在闭包 \(\mathcal{C}^{0}(X;F)\) （\(\mathcal{K}(X;F)\) 在 \(\mathcal{C}^{b}(X;F)\) 中的）上的限制，即在无穷远点趋于 0 的连续函数所成的空间上的限制（Ch. III, §1, No.~2, 命题 3），因此是积分到 \(\mathcal{C}^{0}(X;F)\) 上的连续延拓。
% semantic-fragments: legacy-input-seam
\relax{}
